\documentclass[11pt]{article}

\usepackage[margin=1in]{geometry}
\usepackage{amsmath, amssymb}
\usepackage{booktabs}
\usepackage{listings}
\usepackage{xcolor}
\usepackage{hyperref}

\definecolor{codegray}{gray}{0.95}
\definecolor{commentgreen}{rgb}{0.1,0.45,0.1}
\definecolor{stringred}{rgb}{0.65,0.1,0.1}

\lstset{
  language=R,
  basicstyle=\ttfamily\footnotesize,
  commentstyle=\color{commentgreen}\itshape,
  stringstyle=\color{stringred},
  keywordstyle=\color{blue},
  backgroundcolor=\color{codegray},
  breaklines=true,
  breakatwhitespace=false,
  columns=fullflexible,
  keepspaces=true,
  frame=single,
  rulecolor=\color{black!25},
  showstringspaces=false,
  aboveskip=0.8em,
  belowskip=0.8em
}

% Placeholder for numbers to be copied from the R console
\newcommand{\fillin}[1]{\textbf{[#1]}}

\title{Problem Set 2\\\large Applied Statistics / Quantitative Methods 1}
\author{Mackenzie}
\date{October 2026}

\begin{document}
\maketitle

%==========================================================
\section*{Question 1: Bribery and Driver Class}
%==========================================================

The data cross-tabulate driver class (upper, lower) against the outcome of a traffic encounter (not stopped, bribe requested, stopped/warning), for a total of $n = 42$ drivers.

\subsection*{1(a) Chi-square statistic}

\begin{lstlisting}
##Make the contingency table 
con_tab <- matrix(c(14, 6, 7,
                7, 7, 1),
              nrow = 2, byrow = TRUE,
              dimnames = list(Class   = c("Upper", "Lower"),
                              Outcome = c("Not stopped", "Bribe requested", "Stopped/warning")))
print(con_tab)


# Step 1: totals of each row, column and the overall test
row_tot <- rowSums(con_tab)
col_tot <- colSums(con_tab)
n       <- sum(con_tab)

print(row_tot)
print(col_tot)
print(n)

\end{lstlisting}

The observed counts, with row and column totals, are:

\begin{center}
\begin{tabular}{lcccc}
\toprule
 & Not stopped & Bribe requested & Stopped/warning & Total \\
\midrule
Upper & 14 & 6 & 7 & 27 \\
Lower & 7 & 7 & 1 & 15 \\
\midrule
Total & 21 & 13 & 8 & 42 \\
\bottomrule
\end{tabular}
\end{center}

\begin{lstlisting}
# Step 2: expected counts = (row total * column total) / grand total
# multiplies every row total by every column total, giving the full grid at once
##this is the distribution according to expected counts under independence
expected <- (row_tot %*% t(col_tot)) / n
print(expected)

\end{lstlisting}

Under the null hypothesis of independence between class and outcome, the expected count in each cell is
\[
E_{ij} = \frac{(\text{row total}_i)(\text{column total}_j)}{n}.
\]
This gives:

\begin{center}
\begin{tabular}{lccc}
\toprule
 & Not stopped & Bribe requested & Stopped/warning \\
\midrule
Upper & 13.50 & 8.36 & 5.14 \\
Lower & 7.50 & 4.64 & 2.86 \\
\bottomrule
\end{tabular}
\end{center}

\begin{lstlisting}
# Step 3: chi-square statistic = sum of ((O)bserved - (E)xpected ^2 / E over all cells)
chi_sq <- sum((con_tab - expected)^2 / expected)
print(chi_sq)


#alternativly can also do this : 
result <- chisq.test(con_tab)
print(result)
#extract the expected counts directly
result$expected
\end{lstlisting}

The chi-square statistic is
\[
\chi^2 = \sum_{i,j} \frac{(O_{ij} - E_{ij})^2}{E_{ij}} \approx 3.79.
\]
The built-in \texttt{chisq.test()} returns the same value, which confirms the manual calculation.


\subsection*{1(b) Significance at $\alpha = 0.10$}

The table is $2 \times 3$, so the degrees of freedom are $(2-1)(3-1) = 2$. The p-value is
\[
p = P(\chi^2_2 \geq 3.79) \approx 0.150.
\]
At $\alpha = 0.10$ the critical value is $\chi^2_{2,\,0.10} \approx 4.605$. Since $3.79 < 4.605$ (equivalently, $p = 0.150 > 0.10$), we \textbf{fail to reject} the null hypothesis. This result means that there is not enough statistical evidence at the 10\% significance level to suggest that the variables are related.

\begin{lstlisting}
# Degrees of freedom and p-value
df <- (nrow(con_tab) - 1) * (ncol(con_tab) - 1)
p_val  <- pchisq(chi_sq, df = df, lower.tail = FALSE)

print(df)
print(p_val)


#so the question asks what happens when the alpha is equal to 0.1 or 10%
alpha <- 0.1

if (p_val <= alpha){
  print("results are statistically significant. Reject the null hypothesis")
}else{
  print("results are not statistically significant. Fail to reject the null hypothesis")
}

\end{lstlisting}

\subsection*{1(c) Standardized residuals}

For each cell the standard error is
\[
\mathrm{se}_{ij} = \sqrt{E_{ij}\left(1 - \frac{\text{row total}_i}{n}\right)\left(1 - \frac{\text{column total}_j}{n}\right)},
\]
and the standardized residual is $(O_{ij} - E_{ij})/\mathrm{se}_{ij}$. Each cell is shown below as observed (expected) [standardized residual]:

\begin{center}
\begin{tabular}{lccc}
\toprule
 & Not stopped & Bribe requested & Stopped/warning \\
\midrule
Upper & 14 (13.50) [0.32] & 6 (8.36) [$-1.64$] & 7 (5.14) [1.52] \\
Lower & 7 (7.50) [$-0.32$] & 7 (4.64) [1.64] & 1 (2.86) [$-1.52$] \\
\bottomrule
\end{tabular}
\end{center}

\begin{lstlisting}
#standard error for each cell 
row_prop <- row_tot / n
col_prop <- col_tot / n
se <- sqrt(expected * ((1 - row_prop) %*% t(1 - col_prop)))
print(se)

#standardized residuals: (O - E) / se
st_re <- (con_tab - expected) / se
print(round(st_re, 3))
\end{lstlisting}

\subsection*{1(d) Interpretation}

Under independence, the standardized residuals are approximately standard normal. At $\alpha = 0.10$ (two-tailed, $0.05$ in each tail) the cutoff is $\pm 1.645$.

\begin{itemize}
  \item \textbf{Bribe requested:} the residuals come out at +1.64 for lower class drivers and -1.64 for upper class drivers. So just under the alpha (+/- 1.645)- meaning that lower class drivers get asked for bribes more than independence expected. However, still before the alpha cutoff so it is not statistically significant
  \item \textbf{Stopped/warning:} upper-class drivers were stopped or warned somewhat more than expected ($+1.52$), again below the cutoff.
  \item \textbf{Not stopped:} both residuals are small ($\pm 0.32$), so the two classes were not stopped at roughly the expected rates.
\end{itemize}

No cell reaches significance, which is consistent with the overall test in 1(b). The residuals add information that the overall test cannot: they show \emph{which cells} are furthest from independence and in which direction.

Because the two classes have different sizes, the row percentages are easier to compare than raw counts:

\begin{center}
\begin{tabular}{lccc}
\toprule
 & Not stopped & Bribe requested & Stopped/warning \\
\midrule
Upper & 51.9\% & 22.2\% & 25.9\% \\
Lower & 46.7\% & 46.7\% & 6.7\% \\
\bottomrule
\end{tabular}
\end{center}

Roughly the same share of each class was not stopped, but lower-class drivers were asked for bribes about twice as often (46.7\% vs.\ 22.2\%) and were rarely stopped or warned (6.7\% vs.\ 25.9\%). With only 42 observations, though, these differences are not statistically distinguishable from chance.

\begin{lstlisting}
tab_out <- matrix(sprintf("%.0f (%.2f) [%.2f]", con_tab, expected, st_re),
                  nrow = nrow(con_tab), dimnames = dimnames(con_tab))
noquote(tab_out)

round(100 * prop.table(con_tab, margin = 1), 1)

\end{lstlisting}

%==========================================================
\section*{Question 2: Women Leaders and Drinking-Water Facilities}
%==========================================================

\subsection*{2(a) Hypotheses}

Do villages with a female-reserved leader get more new or repaired drinking-water facilities? The effect of reservation is the slope $\beta_1$ in
\[
\text{water}_i = \beta_0 + \beta_1\,\text{reserved}_i + \varepsilon_i.
\]

   \begin{itemize}
     \item $H_0: \beta_1 = 0$. Reserving the GP head position for women has no effect on the number of new or repaired drinking-water facilities.
     \item $H_1: \beta_1 \neq 0$. Reserving the GP head position for women has an (unknown: either positive or negative) effect on the number of new or repaired drinking-water facilities.
   \end{itemize}


\subsection*{2(b) Bivariate Regression of Hypotheses}

Let $x_i$ be \texttt{reserved}, $y_i$ be \texttt{water}, and $n$ the number of villages. The least-squares estimates are
\[
\hat\beta_1 = \frac{\sum_i (x_i - \bar x)(y_i - \bar y)}{\sum_i (x_i - \bar x)^2},
\qquad
\hat\beta_0 = \bar y - \hat\beta_1 \bar x.
\]
The residual variance uses $n-2$ degrees of freedom because two parameters were estimated from the data:
\[
\hat\sigma^2 = \frac{\sum_i \hat\varepsilon_i^2}{n-2},
\qquad
\mathrm{se}(\hat\beta_1) = \sqrt{\frac{\hat\sigma^2}{\sum_i (x_i - \bar x)^2}}.
\]
The test statistic is $t = \hat\beta_1 / \mathrm{se}(\hat\beta_1)$, compared with the critical value $t_{n-2,\,0.975}$, and the 95\% confidence interval is $\hat\beta_1 \pm t_{n-2,\,0.975}\,\mathrm{se}(\hat\beta_1)$.

\begin{lstlisting}
women <- read.csv("https://raw.githubusercontent.com/kosukeimai/qss/master/PREDICTION/women.csv")

#n = observations, x = reserved, y = water
x <- women$reserved
y <- women$water
nn <- length(y)

print(x)
print(y)
print(nn)

#Compute the means x_bar and y_bar.
x_bar <- mean(x)
y_bar <- mean(y)

print(x_bar)
print(y_bar)

#Intercept: b0_hat = y_bar - b1_hat*x_bar
#Slope: b1_hat = sum(x_i - x_bar)(y_i - y_bar) / sum(x_i - x_bar)^2
Sxy <- sum((x - x_bar) * (y - y_bar))
Sxx <- sum((x - x_bar)^2)
b1  <- Sxy / Sxx
b0  <- y_bar - b1 * x_bar

print(Sxy)
print(Sxx)
print(b1)
print(b0)

fitted_values <- b0 + b1 * x                    
residuals <- y - fitted_values           
sigma2 <- sum(residuals^2) / (nn - 2)     
se_b1  <- sqrt(sigma2 / Sxx)              

print(fitted_values)
print(residuals)
print(sigma2)
print(se_b1)

\end{lstlisting}

\begin{itemize}
  \item \textbf{Fitted Values ($\hat\beta_0$ + $\hat\beta_1$ * x):} prediction for each village 
  \item \textbf{Residuals:} actual minus predicted
  \item \textbf{Residuals Variance (sigma2):} measure of how noisy the data is and nn - 2 instead of nn because we estimated two things (b0 and b1) from the data, which uses up two "degrees of freedom."
  \item \textbf{Standard Error (se\_b1):} is the standard deviation of the estimate across different samples- turns the noise into uncertainty range
\end{itemize}

\begin{lstlisting}

#t-statistic, critical value, two-tailed p-value
t_stat <- b1 / se_b1
t_crit <- qt(0.975, df = nn - 2)
p_value  <- 2 * pt(-abs(t_stat), df = nn - 2)

print(t_stat)
print(t_crit)
print(p_value)

# 95\% confidence interval
ci <- c(b1 - t_crit * se_b1, b1 + t_crit * se_b1)
print(ci)

results <- c(n = nn, b0 = b0, b1 = b1, se = se_b1,
             t = t_stat, t_crit = t_crit, p = p_value)
print(results)


#alternatively, can do this too 
summary(lm(water ~ reserved, data = women))

\end{lstlisting}

The $p$-value answers: If the null were true and reservation had no effect ---> how likely is it that I'd see an estimate this far from zero?
If $p < 0.05$, that's unlikely, so you reject $H_0$. Estimated regression coefficient is large relative to its standard error, meaning the relationship is statistically significant
If $p \geq 0.05$, you fail to reject $H_0$. You don't ``accept'' it, you just lack enough evidence against it.

\begin{lstlisting}

if (abs(t_stat) > t_crit) {
  print("Reject the null hypothesises (H0: b1 = 0), as the t-stat was larger than the critical value")
  }else{
  print("Fail to reject H0- not enough statisticsal evidence  against it")
}
\end{lstlisting}

\begin{center}
\begin{tabular}{lc}
\toprule
Quantity & Value \\
\midrule
Number of villages $n$ & 322 \\
Intercept $\hat\beta_0$ & 14.73832 \\
Slope $\hat\beta_1$ & 9.25 \\
Standard error $\mathrm{se}(\hat\beta_1)$ & 3.947746 \\
$t$-statistic & \fillin{t} \\
Critical value $t_{n-2,\,0.975}$ & 2.343723 \\
Two-tailed $p$-value & 0.0197 \\
95\% CI for $\beta_1$ & 1.485608 17.019238 \\
\bottomrule
\end{tabular}
\end{center}

Since $p = 0.0197 < 0.05$ (equivalently, $|t| > t_{\text{crit}}$), we \textbf{reject} $H_0$. The p-value answers the question: if reservation truly had no effect, how likely would we be to see an estimate this far from zero? Here that is unlikely. Had $p \geq 0.05$, we would only \emph{fail to reject} $H_0$, which is not the same as accepting it. The manual calculation can be checked against \texttt{summary(lm(water $\sim$ reserved, data = women))}.


\subsection*{2(c) Interpret the Coefficient Estimate for Reservation Policy}

The bivariate regression compares two groups of villages: those with a reserved seat ($\text{reserved}=1$) and those without ($\text{reserved}=0$). The regression is a way of making that comparison and testing whether the gap is real or due to chance.

\begin{itemize}
  \item \textbf{Intercept ($\hat\beta_0$):} the average number of new or repaired water facilities in unreserved villages.
  \item \textbf{Slope ($\hat\beta_1$):} the difference between the two group averages (reserved minus unreserved).
  \item \textbf{Standard error:} how much $\hat\beta_1$ would vary across different samples of villages.
  \item \textbf{$t$-statistic:} how many standard errors $\hat\beta_1$ is from zero. A magnitude above roughly 2 means the difference is too large to attribute to chance.
  \item \textbf{$p$-value:} if $H_0$ were true, the probability of an estimate at least this extreme. A small value (here $p < 0.05$) means a difference this large is unlikely to be chance.
\end{itemize}

The estimated coefficient on \texttt{reserved} is 9.25: on average, villages in GPs reserved for women leaders have about 9 more new or repaired drinking-water facilities than villages in unreserved GPs. The p-value of 0.0197 is below 0.05, so we reject the null hypothesis and conclude there is statistically significant evidence that the reservation policy affects the number of drinking-water facilities. Because reservation was randomly assigned, this estimate can be interpreted as the causal effect of reservation.


\end{document}